\section{Objective}
\label{sec:objective}
The primary objective of this thesis is to design, simulate, and evaluate a decentralized \gls{fl} framework for \gls{lorawan}-enabled indoor path loss prediction that overcomes the non-\gls{iid} client drift challenge. Specifically, the thesis evaluates the technical feasibility of deploying a compact 9:8:1 neural network (consisting of exactly 89 parameters) as a \gls{tinyml} model running on the Arduino MKR~WAN~1310. These microcontrollers served as the physical edge nodes (end devices) in the original path loss measurement campaign conducted by Obiri and Van Laerhoven~\cite{obiri2026environment,obiri2025comprehensive}. The research is guided by the research questions summarized in Table~\ref{tab:research-questions}.

\begin{table}[htbp]
\centering
\caption{Research questions guiding this study}
\label{tab:research-questions}
\begin{tabularx}{\textwidth}{lX}
\toprule
\textbf{\#} & \textbf{Research Question} \\
\midrule
RQ1 & Can a federated \gls{tinyml} model achieve comparable path loss prediction accuracy to a centralized baseline trained on the full dataset? \\
\addlinespace
RQ2 & How does non-\gls{iid} data distribution---arising from devices in physically distinct rooms---affect federated model convergence, and to what extent can advanced optimization algorithms (FedProx, SCAFFOLD, and FedYogi) mitigate client drift? \\
\addlinespace
RQ3 & What bandwidth reduction does transmitting model weight updates ($\approx$\,52--104\,B/round) achieve over raw sensor data transmission? \\
\addlinespace
RQ4 & Does the estimated memory footprint and communication payload of the federated \gls{tinyml} model satisfy the physical constraints of the Arduino MKR~WAN~1310, demonstrating feasibility for future real-world deployment? \\
\bottomrule
\end{tabularx}
\end{table}
